% Folder 91, batch 06 — archivists' pages 101–120.
% Pass: Opus 5.5 (claude-opus-5-5), 2026-10-03 — first pass, unchecked against the pages by a human.
%
% Source: archives/batches/91/batch-06.pdf (cover sheet excluded), read from
% the 700-dpi tiles in archives/tiles/91/p101 … p120. Skipped: 101, 118, 120
% (typed Bourbaki pages with nothing of his on them) and 102, 111, 116 (blank).
% The hand is a fast draft: statements and formulas mostly carry, the
% connective prose often does not, and the apparatus says so.
\documentclass[11pt,a4paper]{article}
\input{../preamble/grothendieck}

\folder{91}
\batch{6}
\pages{101}{120}
\foldertitle{Autour de Néron / Greenberg-Néron. Foncteurs Hom (méthodes non-projectives) : notes manuscrites (s.d.), lettre (1967).}
\dating{[à partir de 1964-vers 1970]}
\watermark{Édition de démonstration}

\begin{document}

\section{Topologie constructible}

\page{103}
\note{la page commence par une définition numérotée 1 ; rien n'indique ici une suite d'un argument antérieur au lot}
\textbf{Définition 1.} Soit $X$ un espace topologique. On appelle \emph{topologie constructible} de $X$, \struck{\ill{}} la topologie ayant pour ouverts les réunions quelconques de parties loc.\ constructibles [donc une base d'ouverts formée par les parties loc.\ constructibles]. \uncertain{Partition} \ill{} $=$ \ill{} finie de $X^{\mathrm{cons}}$ $=$ \uncertain{recouvrement} \ill{} par des parties loc.\ constructibles.

\textbf{Prop.~1.} Supposons que \struck{\ill{}} $X$ admette une base \add{$(U_i)_{i\in I}$} d'ouverts \struck{\ill{}} rétrocompacts. Alors
\begin{itemize}
\item[(0)] Les parties constructibles de $X$ forment une base d'ouverts de $X^{\mathrm{cons}}$.
\item[(i)] La topologie constructible de $X$ est plus fine que la topologie initiale, i.e.\ $X^{\mathrm{cons}}\to X$ est continue.
\item[(ii)] \struck{\ill{}} Si $X$ satisfait en plus $T_0$, alors $X^{\mathrm{cons}}$ est séparé [\uncertain{totalement} discontinu] et \uncertain{régulier} \ill{}.
\item[(iii)] Si $E$ est une \struck{partie quelconque de $X$, \ill{} topologique} partie \uncertain{localement} constructible en $X$, munie de la topologie induite $T_E$, alors la topologie constructible \ill{} : $T_E^{\mathrm{cons}}$ est induite par $T^{\mathrm{cons}}$ sur $E$.
\end{itemize}

(0) Une partie loc.\ constructible $E$ est réunion \uncertain{finie} de parties constructibles \uncertain{en} $X$ par $0_{\mathrm{III}}$~9.1.8 [loc]\note{renvoi probable à EGA, chap.~$0_{\mathrm{III}}$, 9.1.8 ; seule la cote « $0_{\mathrm{III}}$ 9.1.8 » est lisible}, \ill{} par les $U_i$ qui sont constructibles.
(i) Résulte \uncertain{du fait} \ill{} \ill{} pour $T^{\mathrm{cons}}$.
Pour (ii), soient $x,y\in X$, il existe \ill{} un ouvert \ill{} contenant $x$, et non $y$ \ill{} : $U\ni x$, $U\not\ni y$. On peut supposer, \uncertain{grâce à} l'hyp., $U$ rétrocompact, donc [loc.] constructible\add{, donc ouvert}, et $X-U=V$ l'est \uncertain{aussi}, donc $U\cap V=\emptyset$, $U\cup V=X$, $U\ni x$, $V\ni y$, $U$, $V$ ouverts pour $T^{\mathrm{cons}}$.
Pour (iii), on note que la topologie induite \ill{} $T^{\mathrm{cons}}$ sur $E$ a une base d'ouverts les $E\cap F$, $F$ \add{loc.} constructible dans $X$, or \ill{} les $G\subset E$ loc.\ constructibles dans $E$, \ill{}

\page{104}
\ill{} \ill{}

\textbf{Lemme.} a) Soit $F$ (loc.) constructible dans $X$, alors $F\cap E$ est (loc.) constructible \uncertain{dans} $E$ \add{($E$ loc.\ const.)}.
b) Soit $F\subset E$, \add{$E$ (loc.) constructible} alors $F$ est (loc.) const.\ dans $E$ ssi il l'est dans $X$.

\marginal{$E$, $f^{-1}(F)\cap U$ ; $f$ ; $X$, $F$}\note{notation griffonnée dans la marge gauche, en face de la démonstration ; lecture et rôle incertains}

Dém. Pour a), on \uncertain{note} que si $F$ est \uncertain{ouvert} rétrocompact en $X$, alors $F\cap E$ est rétrocompact dans $E$, \ill{} \ill{} que $A$ \ill{} quasi-compact \add{$U$} \ill{} $E$ est \ill{} \uncertain{par} un ouvert $V$ de $X$, quasi-compact \ill{} [\ill{} : l'hyp.\ préliminaire de la \uncertain{prop.}] \uncertain{supposons} quasi-compact, d'où $f^{-1}(F)\cap U=f^{-1}(F\cap V)=$ \ill{} \ill{} d'un quasi-compact, \ill{} \ill{} \ill{} \ill{} $E$ est \ill{} rétrocompact en $X$. \struck{\ill{}} \add{$0_{\mathrm{III}}$~9.1.13} \ill{} \ill{} « localement » \ill{} \ill{}, \uncertain{utilisons} le fait que $E$ \uncertain{est} \ill{} \ill{} $E\cap U$ constructible \uncertain{dans} \uncertain{tout} $U$ \ill{}.

Pour b), on suppose $F\subset E$ constructible dans $E$ \add{constr.}, \uncertain{prouvons} \ill{}, que $F$ l'est dans $X$. On \ill{} $F$ \struck{\ill{}} \add{ouvert} dans $E$, rétrocompact en $E$, \uncertain{donc} $F=U\cap \complement V$, $U$ et $V$ rétrocompacts dans $X$ \add{qui sont constructibles} \ill{}

[\uncertain{Or} $E=\bigcup E_i$ ; \uncertain{\ill{}} \ill{} \ill{}, \ill{} \struck{$F\cap E_i$ \ill{}} \ill{} $E$ \ill{} \ill{}, donc $E_i\cap F$ \uncertain{const.} \ill{} $E_i$ \ill{} \ill{}, \ill{} \ill{} \ill{} \ill{} : \ill{} \ill{} \uncertain{qu'ils sont} \uncertain{constructibles} dans $X$, \ill{} \ill{} $F=\bigcup E_i\cap F$. $0_{\mathrm{III}}$~9.1.8 (i)$\Rightarrow$(ii)\note{renvoi probable à EGA $0_{\mathrm{III}}$ 9.1.8 ; un « (ii) » est cerclé juste au-dessous}

\uncertain{Remplaçons} $X$ \ill{} $U$, \ill{} \ill{} \ill{} $E$, \ill{}
\struck{$F$ \ill{} \ill{} \ill{} dans \ill{}} \ill{} \ill{} \ill{} $X=U$ [$E\subset$ \ill{} \uncertain{const.}\ dans $U$, $F$ \ill{}]
\struck{$F=E\cap W$, $W$ \ill{} dans $X$,} \struck{\ill{}} \uncertain{de même} dans $U$, \uncertain{donc} dans $X$]
\note{la fin de la page est très raturée ; plusieurs lignes se chevauchent et la ponctuation n'est pas sûre}

\drawing{dans la marge inférieure gauche, un ovale traversé par un trait oblique}

\page{105}
Dans \ill{} \ill{} \uncertain{supposons} $E=\complement V$, $V$ \struck{\ill{}} ouvert rétrocompact, \struck{et \ill{} \ill{} \ill{} \ill{} $E-F=\complement$ \ill{} \ill{}} \uncertain{Il suffit} \uncertain{de prouver} \add{la \uncertain{\ill{}}} $E-F$ \uncertain{const.}\ en $X$, i.e.\
\[
X-(E-F)=V\cup F \quad \text{rétrocompact dans } X,
\]
\uncertain{ce qui est} \ill{} \ill{} \ill{} $V$ et $F$ le sont \uncertain{évidemment} \ldots
\marginal{$\complement V$ ; $F$}\drawing{dans la marge gauche, un segment oblique gradué de deux tirets, marqué $\complement V$ en haut et $F$ au milieu}

Le cas « local » résulte aussitôt \uncertain{de ce} \uncertain{lemme}.

\textbf{Corollaire 1.} Supposons \uncertain{que} \ill{} \ill{} \uncertain{plus}, \uncertain{toute} partie $U_i$ soit quasi-compacte \struck{\ill{} \ill{}} \uncertain{[pour la topologie]} constructible [donc \ill{} \struck{\ill{} \ill{}} compacte \uncertain{pour} la top.\ constructible, donc \ill{} \ill{} \ill{} \struck{\ill{}} $X$ i.e.\ satisfaisant à la condition : \struck{\ill{}} \ill{} \ill{} séparée $(T_0)$], \struck{$A$} \ill{} \ill{}
\begin{itemize}
\item[(iv)] Les $U_i^{\mathrm{cons}}$ sont quasi-compacts, \struck{et \ill{}} \struck{\ill{} $X$}
\struck{\ill{} \ill{} \ill{} \ill{} \ill{} \ill{} \ill{} \ill{} \ill{} \ill{}} \ill{} \ill{} \ill{} \struck{[\ill{} $X$ \ill{} \ill{} \ill{} \ill{} \ill{}]} \note{deux lignes biffées, avec une accolade ondulée, avant (v)}
\item[(v)] Soit $E$ une partie \struck{loc.\ constructible de} \ill{} \ill{} $X^{\mathrm{cons}}$ [\uncertain{p.\ ex.} une partie loc.\ constructible \add{(\uncertain{fermée} quasi-compacte)} \ill{}] de $X$] dans $X^{\mathrm{cons}}$. \ill{} $E$ \struck{compacte} quasi-compacte dans $X^{\mathrm{cons}}$ $\iff$ $E$ quasi-compact dans $X$ $\iff$ $E\subset$ réunion finie des $U_i$.
\item[(vi)] \uncertain{Si} \ill{} \uncertain{que} $X^{\mathrm{cons}}$ soit loc.\ compact [\uncertain{compact}] il \uncertain{faut et suffit} que $X$ satisfasse $(T_0)$ [\uncertain{compact} $T_0$, et \uncertain{soit} quasi-compact].
\end{itemize}
(iv) Résulte de (i) \uncertain{trivialement}. \uncertain{Évidemment}
\[
E^{\mathrm{cons}} \text{ qu.-cpte} \Longrightarrow E \text{ qu.\ cpct} \quad \text{(i)}
\]
\[
\Longrightarrow E \text{ est contenu dans une réunion finie des } U_i
\]
par définition de qu.\ cpct,\note{dans ces deux lignes, « contenu » et « réunion finie » sont des lectures incertaines ; trois mots illisibles suivent « qu.\ cpct » dans la première}
donc $E$ fermé dans $\bigcup U_i^{\mathrm{cons}}$ \add{$V^{\mathrm{cons}}$ [par la top.\ induite par $X^{\mathrm{cons}}$]}, et les $U_i^{\mathrm{cons}}$ étant quasi-compacts \ill{} \ill{} \uncertain{il s'ensuit} que $E$ l'est, OK.
(vi) En \uncertain{résulte} \uncertain{trivialement}.

\marginal{N.B. Les hypothèses faites impliquent \uncertain{que} \ill{} \ill{} \uncertain{des} \uncertain{points} \ill{} de $X$ \ill{} \ill{} \uncertain{quasi-compacts}, \ill{} \ill{} \uncertain{donc} $(T_0)$ \ill{} \uncertain{quelconques} \ill{}}\note{note écrite en oblique dans la marge gauche, à hauteur des énoncés (iv)--(v) ; lecture très incomplète}

\page{106}
\textbf{Corollaire 2.} \struck{\ill{} les conditions des corollaires} Supposons $X$ \struck{\ill{}} quasi-compact \add{et $(T_0)$}, et que \struck{la partie} $X$ ait une base formée d'ouverts \struck{\ill{}} rétrocompacts, qui soient quasi-compacts pour la topologie constructible. Soit $E$ une partie de $X$. \uncertain{Conditions} équivalentes :
\begin{itemize}
\item[a)] $E$ est fermé pour la top.\ constructible [i.e.\ est \uncertain{intersection} \struck{\ill{}} de parties constructibles]
\item[b)] $E$ est compact pour la top.\ constructible, i.e.\ pour toute famille de parties constructibles $(E_i)_{i\in I}$ de $X$, \uncertain{avec} $E\subset\bigcup_{i\in I}E_i$, il existe une partie \uncertain{finie} $J\subset I$, telle que $E\subset\bigcup_{i\in J}E_i$.
\end{itemize}
\marginal{Notion locale ?? ; \uncertain{Sous}-familles \uncertain{finies} \ill{} ; \uncertain{partie} \uncertain{ensembliste} \uncertain{intersection} \uncertain{de} \uncertain{localement} \ill{} \ill{}}\note{notes obliques dans la marge gauche, en partie barrées ; les « ?? » sont de l'auteur}

\struck{$X$ \ill{}} \uncertain{En effet}, $X^{\mathrm{cons}}$ est compact, donc a)$\iff$b).

\uncertain{On en conclut} \add{\uncertain{donc}} \uncertain{que} \struck{\ill{}} $E$ est une partie \uncertain{ensembliste} \uncertain{fermée} \uncertain{dans} $X$. On \uncertain{peut} \ill{} \uncertain{que} $X$ \ill{} \ill{} \ill{} \ill{} \uncertain{plus} \uncertain{précis} \uncertain{ici} [\uncertain{mais} \uncertain{seulement} \uncertain{fermé} \ill{} \ill{} \ill{} \uncertain{deux} \uncertain{top}] \add{dans $X^{\mathrm{cons}}$}

$E$ fermé [i.e.\ intersection de parties loc.\ constructibles]
\[
\Updownarrow
\]
$E$ loc.\ compact pour la top.\ constructible i.e.\ pour tout \struck{\ill{}} $U_i$, $E\cap U_i$ satisfait la condition b) \uncertain{relativement} \uncertain{à} $U_i$.

\textbf{Corollaire 3.} Sous les conditions du corollaire 2, conditions \uncertain{équivalentes} :
\begin{itemize}
\item[a)] $E$ et $X-E$ sont \uncertain{ensemblistes} \struck{constructibles} \add{ouverts dans $X^{\mathrm{cons}}$}
\item[b)] $E$ est constructible.
\end{itemize}

\page{107}
c)$\Rightarrow$a)$\iff$b) \uncertain{trivial}, \uncertain{prouvons} la réciproque a)$\Rightarrow$c), alors \uncertain{soit} $E$ et $X-E=F$ \uncertain{clos} \ill{} en $X^{\mathrm{cons}}$ ; \uncertain{ce sont} \uncertain{réunions} \struck{\ill{} \ill{}} \uncertain{filtrantes} \uncertain{d'ouverts} constructibles $E_i$ \uncertain{resp.}\ $F_j$, \uncertain{donc} $X$ \uncertain{est} \uncertain{réunion} \uncertain{filtrante} \uncertain{des} $E_i\cup F_j$, donc \uncertain{est} \uncertain{égal} \uncertain{à} l'un des $E_i\cup F_j$, d'\uncertain{où} $E=E_i$, \struck{\ill{}} $F=F_j$ \uncertain{pour} un $i,j$, donc $E$, $F$ constructibles.

\textbf{Proposition 2.} Soit $f:X\to Y$ une application continue d'espaces topologiques, telles que \uncertain{$Y$} \uncertain{admette} \uncertain{une} \uncertain{base} \uncertain{formée} d'ouverts \add{$V$ rétrocompacts en $U$}, $V\subset U\subset Y$, \struck{rétrocompacts dans $Y$,} \uncertain{tels} \uncertain{que} $f^{-1}(V)$ soit rétrocompact dans $f^{-1}(U)$. Alors
\struck{(i) \uncertain{Pour} \uncertain{tout} \ill{} $Y'$ de $Y$, \ill{} $X'=f^{-1}(Y')$, $f':X'\to Y'$ satisfait la \ill{}} \add{\uncertain{hyp.}\ \uncertain{par} $f$.}
\begin{itemize}
\item[(i)] Pour toute partie constructible [loc.\ constructible] $E$ de $Y$, $f^{-1}(E)$ est une partie constructible [loc.\ constructible].
\item[(ii)] $f$ est \uncertain{continue} \uncertain{de} $X^{\mathrm{cons}}$ \uncertain{dans} $Y^{\mathrm{cons}}$, i.e.\ l'image inverse d'une partie \add{\uncertain{loc.}} \uncertain{ensembliste}
\end{itemize}
de $Y$ est une partie \add{\uncertain{loc.}} \uncertain{ensembliste} de $X$.
\marginal{hyp.\ trop \uncertain{restrictive} ; \ill{} \ill{} \uncertain{préalables} \ill{}}\note{note oblique de l'auteur dans la marge gauche, en face de l'hypothèse de la proposition 2}
\marginal{$Y'\subset Y$}

(i) est \uncertain{immédiat} \uncertain{pour} « constructible ». \uncertain{Pour} loc.\ constructible, \ill{} \ill{} \ill{} \ill{} $Y=\bigcup Z_i$, \uncertain{les} $E_i=E\cap Z_i$ constructibles \add{\uncertain{dans} $Z_i$}, \uncertain{\ill{}} $Z_i$ \uncertain{ouvert}, $f_i:X_i\to Z_i$ \uncertain{induit}, $X_i=f^{-1}(Z_i)$, [$U$ \uncertain{et} $\complement V$ \ill{} \ill{}] \ill{} \ill{} \uncertain{que} $f_i^{-1}(E_i)$ \uncertain{est} constructible \uncertain{dans} $X_i$ ; \ill{} \ill{} \ill{} \uncertain{que} $X_i\to Z_i$ satisfait l'hypp.\ de $f:X\to Y$, \struck{\ill{}} OK.
(ii) Résulte \uncertain{trivialement} \uncertain{de} (i).
\marginal{$f'$}

\page{108}
\note{la première moitié de la page (corollaire 1 et sa démonstration) est barrée de deux longs traits obliques}
\struck{\textbf{Corollaire 1.} \uncertain{Supposons} que \ill{} $X$ et $Y$, \ill{} \uncertain{admettent} une \ill{} \ill{} \uncertain{d'ouverts} \add{$X_i$ \uncertain{resp.}\ $Y_j$} rétrocompacts quasi-compacts. \ill{} \ill{} \ill{} \ill{} \uncertain{Mais} l'hyp.\ préliminaire de la \uncertain{prop.}\ \uncertain{signifie} \ill{} \ill{} \uncertain{que} \uncertain{pour} \uncertain{tout} \ill{} $Y_j$, $f^{-1}(Y_j)$ est rétrocompact en $X$, \ill{} \ill{} \ill{} \uncertain{tout} \struck{ouvert} \ill{} $Y'$ \add{[\uncertain{rétrocompact}]} \uncertain{quasi-compact} de $Y$, $X'=f^{-1}(Y')$ est \struck{\ill{}} rétrocompact dans $X$, \ill{} \ill{} \ill{} les $X_i\cap f^{-1}(Y_j)$ sont qu.\ cpcts.}

[\textbf{Prop.} Soit $X$ un espace topologique, \uncertain{ayant} une base d'ouverts $U_i$ qui sont quasi-compacts. Alors \struck{\uncertain{pour} \uncertain{toute} partie $E$ de $X$, \ill{} \ill{}} \uncertain{pour} \uncertain{une} application \uncertain{continue} $f:X'\to X$, il \uncertain{revient} \ill{} \ill{} \ill{} \ill{} \ill{} \uncertain{que} $f$ \uncertain{est} qu.\ cpct, i.e.\ que les $f^{-1}(U_i)$ sont q.\ cpcts. Dans \uncertain{ce cas} $X'\subset X$ \uncertain{sous-espace} \uncertain{top.}, $X'$ \uncertain{est} \uncertain{rétrocompact} en $X$ \uncertain{ssi} les $X'\cap U_i$ sont qu.\ cpcts.

\textbf{Cor.} \struck{Soit} \struck{Supposons} \add{\uncertain{Pour} \uncertain{que}} \struck{de plus} les $U_i$ rétrocompacts \add{(il \uncertain{suffit} \uncertain{que})} dans $X$, \struck{\ill{}} \emph{toutes} parties quasi-compactes ouvertes de $X$ soient rétrocompactes en $X$.

\textbf{Cor.} S'il en est ainsi, pour une application \uncertain{continue} $f:X'\to X$, il \uncertain{revient} au \uncertain{même} de \uncertain{dire} :
\begin{itemize}
\item[(a)] \uncertain{que} \uncertain{pour} tout \uncertain{ouvert} $U\subset X$ \add{\ill{}} quasi-compact, $f^{-1}(U)$ est rétrocompact ;
\item[(b)] Pour tout $U_i$, $f^{-1}(U_i)$ est rétrocompact \uncertain{dans} $X'$ ;
\item[(c)] [Si $X'$ \uncertain{a} une base d'ouverts quasicompacts $U'_j$] les $f^{-1}(U_i)\cap U'_j$ sont \struck{\ill{}} quasi-cpcts.]
\end{itemize}

\marginal{$U\subset X$ \uncertain{ouvert} \struck{rétrocompact} \ill{} ; \uncertain{si} $U'\subset X'$ \ill{} q.\ cpct}

b)$\Rightarrow$a), $E=U$, $f^{-1}(U)\cap V'$ $\subset f^{-1}(f(V'))$, \struck{$=f^{-1}(U\cap$ \ill{}$\cap$} \uncertain{et} $f(V')$ \uncertain{est} \uncertain{q.\ cpct}, \uncertain{est} \uncertain{contenu} \uncertain{dans} une \uncertain{réunion} \uncertain{finie} \uncertain{des} $U_i$, $i\in J$, donc $f^{-1}(U)\cap V'=f^{-1}(\bigcup U\cap U_i)\cap V'$ \uncertain{et} \uncertain{chaque} $U\cap U_i$ \uncertain{est} \add{$i\in J$} \uncertain{réunion} \uncertain{finie} \uncertain{d'un} $U_j$ \ldots
\note{ce calcul, écrit dans la marge gauche sous une note entourée, se poursuit vers le bas de la page et s'interrompt}

\struck{\textbf{Cor.} Si $f:X'\to X$ \ill{} \uncertain{quasi-compacts} \ill{} \ill{} \ill{} \ill{} \ill{}}\note{note oblique dans la marge gauche, raturée par un trait ondulé}

\page{109}
\textbf{Corollaire 1.} Supposons que $X$, \add{$Y$} \struck{\ill{} \ill{}} \uncertain{soient} des espaces $(T_0)$, \struck{\uncertain{schématiques} \uncertain{ou}} admettant des bases d'ouverts formées d'\struck{\uncertain{quasi-compacts}} \uncertain{ouverts} rétrocompacts et quasi-compacts pour la top.\ constructible [donc compacts pour la \uncertain{\ill{}}]. \struck{\uncertain{Conditions}} \add{Conditions} \uncertain{équivalentes} :
\begin{itemize}
\item[a)] $f:X\to Y$ \struck{soit} \struck{une \uncertain{telle} \uncertain{que}} \add{est} quasi-compacte
\item[b)] \struck{Alors} \add{(b)} $f^{\mathrm{cons}}$ est \uncertain{une} \struck{\ill{}} \struck{[$f$ \uncertain{compact} \ill{} $X$]} application \uncertain{propre} d'espaces loc.\ cpcts, \struck{\ill{} \uncertain{partie} \ill{}} \struck{i.e.\ $f$ \ill{} \ill{}}
\struck{[\ill{} \ill{} \uncertain{pour} \uncertain{partie} \ill{}]}
\struck{\ill{} \ill{} \ill{} \ill{} \ill{}} \note{lignes biffées et entourées en milieu de page}
\end{itemize}
[N.B. $f$ \uncertain{est} \uncertain{propre} \uncertain{ssi} \uncertain{l'image} \uncertain{réciproque} \uncertain{d'une} partie \uncertain{compacte} \uncertain{est} \uncertain{compacte}.] \struck{\uncertain{les} parties constructibles \uncertain{en} parties \add{\uncertain{\ill{}}} \ill{} \ill{} \ill{} \ill{} \uncertain{fermées} $f^{\mathrm{cons}}$ \ill{}, i.e.\ \uncertain{une} \uncertain{transformée} \ill{}}

\marginal{c) $f$ transforme \uncertain{parties} \uncertain{ensemblistes} \uncertain{en} \uncertain{parties} \uncertain{ensemblistes}, et les \uncertain{fibres} de $f$ \uncertain{sont} quasi-compactes [\uncertain{cf.} Bourbaki, \ill{} I, §~10, 2, Prop.~1]}\note{note oblique de l'auteur dans la marge gauche, séparée du texte par un trait courbe}

\textbf{Corollaire 2.} \uncertain{Soient} \uncertain{les} \uncertain{conditions} \uncertain{précédentes} \add{\uncertain{\ill{}} \uncertain{vérifiées},} supposons $Y$ \uncertain{et} $X$ \uncertain{quasi-cpcts}, \uncertain{de} $X^{\mathrm{cons}}$, $Y^{\mathrm{cons}}$ \uncertain{compacts}. Conditions équivalentes :
\begin{itemize}
\item[(a)] $f$ transforme \uncertain{les} parties constructibles en \uncertain{parties} \uncertain{\ill{}} [i.e.\ transforme « \ill{} \ill{} » en « \ill{} \ill{} »]
\item[(b)] $f$ transforme \uncertain{les} parties \uncertain{ensemblistes} en \uncertain{\ill{}} \uncertain{parties} constructibles [\uncertain{\ill{}} « \ill{} \ill{} » \uncertain{en} « \ill{} »]
\item[(c)] $f^{\mathrm{cons}}$ est une application ouverte. \struck{\ill{} \ill{} \ill{}}
\end{itemize}
Les conditions sont \uncertain{vérifiées} si $f$ est \emph{bijective}.

\textbf{Cor.~3.} Sous les conditions du \uncertain{Cor.~1} \uncertain{ou} \uncertain{supposons}, alors \struck{(\ill{}) \ill{} \ill{}} (ii) \uncertain{pour} \uncertain{toute} partie \add{$E$} \uncertain{de} $f(X)$, \ill{} \uncertain{dont} $f^{-1}(E)$ \uncertain{l'est}\ldots{} (iii) Si $f$ est \uncertain{surjective}, une \uncertain{partie} \ill{} \ill{} \ill{} constructible \uncertain{ssi} $f^{-1}(Y)$ l'\uncertain{est}.
\note{la fin du corollaire 3, serrée au bas de la page et en partie raturée, n'est lue que par fragments}

\page{110}
\uncertain{Examiner} \ill{} \uncertain{plus} \uncertain{de} \uncertain{près} \uncertain{les} parties \uncertain{ensemblistes} \add{(et loc.\ constructibles)} de $X$, \uncertain{quand} \ill{} \ill{} suppose pas $X$ qu.\ cpct\ldots

\textbf{Th.} Soit $X$ un schéma affine. Alors $X^{\mathrm{cons}}$ est compact.

\textbf{Corollaire 1.} Soit $X$ un préschéma. \struck{\ill{}} Alors $X$ \add{qu.\ séparé} est $T_0$, \uncertain{admet} une base formée d'ouverts rétrocompacts et compacts pour la topologie constructible. En particulier, $X^{\mathrm{cons}}$ est loc.\ compact.

\textbf{Corollaire 2.} Supposons \struck{\uncertain{de plus}} $X$ [quasi-compact et quasi-séparé] \struck{\uncertain{quasi-noethérien}}. \struck{Alors} Alors
\begin{itemize}
\item[(i)] $X^{\mathrm{cons}}$ est compact.
\item[(ii)] Les parties constructibles de $X$ sont les parties \uncertain{ouvertes} et fermées de $X^{\mathrm{cons}}$.
\item[(iii)] Les parties \struck{\ill{}} \uncertain{ensemblistes} de $X$ \uncertain{sont} les parties \add{$E$} \uncertain{telles} \uncertain{qu'il} \uncertain{existe} un \uncertain{morphisme} \uncertain{quasi-compact} \add{\uncertain{de}} $f:X'\to X$, \uncertain{d'image} $E$. \struck{[N.B.\ \ill{}]} \uncertain{On} \uncertain{peut} \uncertain{prendre} \struck{$X$} $X'\to X$ \uncertain{un} \uncertain{monomorphisme}\ldots
\item[(iv)] \uncertain{Si} $E$ et $X-E$ \uncertain{sont} qu.\ cpct\ldots, \uncertain{ils} \uncertain{sont} constructibles\ldots
\end{itemize}

\section{Applications constructibles et espaces schématiques}

\page{112}
Soient $X$, $Y$ deux préschémas. \emph{Application \uncertain{géométrique}} de $E\subset X$ \uncertain{dans} $Y$ : \uncertain{morphisme} \uncertain{des} \add{\uncertain{\ill{}} \uncertain{supposés} $E$ \uncertain{loc.}\ constructibles} $\mathrm{Dis}(E,X)=\coprod_{x\in E}\mathrm{Spec}\,k(x)$, \uncertain{dans} $\mathrm{Dis}(Y)$. [Une application \uncertain{géométrique} est dite \emph{constructible} si \uncertain{tout} $x\in X$ a un voisinage \add{\uncertain{\ill{}}} \uncertain{ouvert} $U$, et une partition de $U\cap E$ en parties loc.\ \uncertain{fermées} $Z_i$, telle que $f\,|\,Z_i$ \uncertain{soit} un \emph{morphisme} $Z_i\to Y$. \struck{Il \uncertain{revient} \uncertain{aussi} \ill{} \ill{}} \struck{\ill{} \uncertain{dire} que}

Ex : une appl.\ \uncertain{qui} \uncertain{est} \ill{} \uncertain{morphisme}. Si \ill{}\ldots{} \uncertain{ensembles} \ill{} \ill{} \uncertain{définie} \uncertain{par} un \ill{} \ill{} \ill{} \uncertain{partie} \uncertain{loc.}

$E=\bigcup_{i\in I}E_i$ famille \ill{} \ill{} $f:E\to Y$ \ill{} \uncertain{est} \ill{} \uncertain{constructible}, alors \uncertain{pour} tout $i$, $f\,|\,E_i$ l'est ; \uncertain{la} \uncertain{réciproque} \add{\uncertain{\ill{}} \ill{} \uncertain{est} \uncertain{vraie} \ill{} \ill{} \uncertain{si} \ill{} \ill{} \uncertain{fini}} \struck{\ill{}} \uncertain{constituant} une \emph{partition} \uncertain{si} $E$ \uncertain{est} \uncertain{quasi-compact} \add{\ill{} \uncertain{partition} \uncertain{en} \uncertain{parties} $Z_i$ constructibles} \uncertain{possible} de $E$ \uncertain{en} \ill{} loc.\ \uncertain{fermées}, telles que les $f\,|\,Z_i$ soient des morphismes.

\uncertain{La} \uncertain{définition} \uncertain{de} $\mathrm{App\,cons}((X,E),(Y,F))$, \uncertain{\ill{}}.

\add{(Cons Aff)}\note{ajout de l'auteur au-dessus de la ligne, rattaché par un trait}
\uncertain{Catégorie} \uncertain{des} \uncertain{ensembles} \uncertain{constructibles} \uncertain{affines} : \uncertain{objets} \uncertain{les} \struck{Affines} \ill{} \ill{} \ill{} \ill{}, \uncertain{\ill{}} \uncertain{les} schémas \uncertain{affines} (\uncertain{ou} quasi-compacts \uncertain{\ill{}}) et morphismes les \uncertain{applications} constructibles. Si $X$ \struck{$E$} est une partie loc.\ \struck{\ill{}} \uncertain{préschéma} \ill{} \ill{} constructible d'un préschéma $X$, \uncertain{on} \uncertain{lui} \uncertain{associe} un \uncertain{foncteur}
\[
\mathrm{Cons}(X,E) : (\mathrm{Cons\,Aff})^{\circ}\longrightarrow(\mathrm{Ens})
\]
\note{la page s'arrête sur cette définition ; la suite manque sur ce feuillet}

\page{113}
\uncertain{Posons}
\[
\mathrm{Cons}(X,E)(Y)=\mathrm{App\,cons}(Y,(X,E)).
\]
\struck{\ill{} \ill{}} Si \uncertain{\ill{}}\ldots{} $(X,E)$ \uncertain{et} $(X',E')$ ; \uncertain{on} \uncertain{a}\ldots{}
\[
\mathrm{App\,cons}((X,E),(X',E'))\xrightarrow{\ \sim\ }\mathrm{Hom}(\mathrm{Cons}(X,E),\mathrm{Cons}(X',E'))
\]

\uncertain{On} \uncertain{obtient} \uncertain{ainsi} \uncertain{des} foncteurs pleinement fidèles

\begin{tikzcd}[column sep=small, row sep=large, nodes={font=\scriptsize}]
(\mathrm{Cons\,Aff}) \arrow[r] \arrow[rr, bend right=25, "\text{canonique}"'] \arrow[rrr, bend right=35, "\mathrm{can}"'] & (\mathrm{Complexes}(X,E)) \arrow[r] & \mathrm{Ind}(\mathrm{Cons\,Aff})' \arrow[r] \arrow[r, bend right=40, "\mathrm{can}"'] & \underline{\mathrm{Hom}}((\mathrm{Cons\,Aff})^{\circ},\mathrm{Ens})
\end{tikzcd}

\note{sous $(\mathrm{Complexes}(X,E))$ : « $E$ loc.\ constr.\ dans $X$ \uncertain{quasi-sép.} » ; sous $\mathrm{Ind}(\mathrm{Cons\,Aff})'$ : « systèmes inductifs \uncertain{\ill{}}\ldots ». Au-dessus de la flèche $(\mathrm{Cons\,Aff})\to(\mathrm{Complexes}(X,E))$, une note de l'auteur, en partie lisible : « \uncertain{équivalence} \uncertain{avec} \uncertain{les} \uncertain{ens.}\ \uncertain{schématiques} \uncertain{quasi-compacts}\ldots ». Les trois arcs courbes sont redessinés d'après leurs extrémités apparentes ; le départ de l'arc court marqué « can » (sous $\mathrm{Ind}$ ou sous $\mathrm{Complexes}$) est incertain}
\note{le reste de la page est blanc}

\page{114}
Soit \struck{$A$} $A$ un anneau commutatif, $X$ le schéma affine qu'il définit,
\[
A^{\mathrm{cons}}=\varinjlim_{\mathfrak{P}} A_{\mathfrak{P}}
\]
$\mathfrak{P}$ parcourant les « partitions \uncertain{ordinaires} » de $X$ en ens.\ constructibles loc.\ fermés, affines, et pour une telle $\mathfrak{P}=(X_i)_{i\in I}$,
\[
A_{\mathfrak{P}} \text{ désigne } A\Bigl(\coprod X_{i\,\mathrm{red}}\Bigr)=\prod A(X_{i\,\mathrm{red}}).
\]
\note{après « désigne », un mot biffé illisible. Le $\mathfrak{P}$ est une lettre gothique majuscule, sans doute un P ; lecture de la lettre incertaine}
Les homom.\ de transition sont injectifs. Considérons
\[
\widetilde{X}=\mathrm{Spec}(A^{\mathrm{cons}}).
\]
On a
\[
\mathrm{Top}(\widetilde{X})=\varprojlim_{\mathfrak{P}}\mathrm{Top}\,\mathrm{Spec}(A_{\mathfrak{P}}).
\]
En particulier
\[
\mathrm{Ens}(\widetilde{X})=\varprojlim_{\mathfrak{P}}\mathrm{Ens}\underbrace{\mathrm{Spec}\,A_{\mathfrak{P}}}_{\coprod X_i}
\]
et comme les \struck{\uncertain{homéomorphismes}} \add{bijections} \uncertain{morphismes} de transition sont des \struck{\ill{}} \add{bijections}, on trouve
\[
\mathrm{Ens}\,\widetilde{X}\simeq\mathrm{Ens}\,X.
\]
Cette bijection est compatible avec les topologies de $\widetilde{X}$ et de $X^{\mathrm{cons}}$, comme on voit \uncertain{aisément} sur la description d'une base d'ouverts de $\widetilde{X}$\ldots

\page{115}
\note{toute la page est barrée de grands traits obliques croisés ; un signe $\emptyset$ cerclé et barré se trouve dans la marge gauche, en face du N.B.}
Soit $X$ un espace topologique, \struck{dont chaque point $x$ a un \ill{}} \add{admettant une base d'ouverts} \ill{}.
\marginal{\uncertain{\ill{}}}\note{mot isolé de l'auteur en marge gauche, en face de la première ligne}
\struck{\ill{} \ill{} \ill{} \ill{} \ill{}} \uncertain{que} $U$ \uncertain{ayant} les propriétés \uncertain{suivantes} :
\begin{itemize}
\item[(i)] $U$ rétrocompact dans $X$, et $U$ quasi-compact
\item[(ii)] Toute \struck{\ill{}} famille $(E_i)_{i\in I}$ de parties \struck{\ill{}} constructibles de $U$, recouvrant $U$, il y a une sous-famille finie recouvrant $U$.
\end{itemize}

Pour exprimer (ii), \uncertain{introduisons} [dans \uncertain{tout} espace topologique $U$] la topologie \uncertain{dont} les \uncertain{ouverts} \uncertain{sont} les \add{réunions quelconques des} \add{\uncertain{parties}} loc.\ constructibles [\uncertain{c'est} bien \uncertain{stable} par \uncertain{intersections} finies et \uncertain{réunions} \uncertain{quelc.}] \uncertain{dans} \uncertain{la} \uncertain{topologie} \uncertain{initiale}, \uncertain{appelée} \uncertain{la} \uncertain{topologie} \uncertain{des} \uncertain{parties} \uncertain{constructibles} ($U$ est quasi-compact, et \uncertain{donc} \uncertain{on} \uncertain{a} \uncertain{une} base d'ouverts \uncertain{formée} \uncertain{de} \uncertain{parties} \uncertain{constructibles}. \uncertain{Lorsque} \ill{} \struck{\ill{} \ill{}} \uncertain{admet} \ill{} \uncertain{ouverts} rétrocompacts, « loc.\ constructible » et \uncertain{d'ailleurs} \uncertain{toujours} \uncertain{ouverts} « constructibles » ; \uncertain{lorsque} (ii), \uncertain{les} \struck{$U$} \uncertain{ouverts} rétrocompacts \uncertain{dans} \uncertain{\ill{}} \uncertain{X}) \uncertain{forment} \uncertain{donc} \uncertain{une} \ill{} de la \struck{\ill{}} \ill{}

\uncertain{au} \uncertain{précédent}. N.B.\ \add{\uncertain{\ill{}}} \uncertain{Il} \uncertain{en} \uncertain{est} \struck{\ill{}} \uncertain{ainsi} \struck{\ill{}} [\uncertain{c'est} \uncertain{le} \uncertain{cas} \uncertain{des} \uncertain{préschémas}] \uncertain{dans}
l'\uncertain{axiome} $T_0$, alors \uncertain{une} \uncertain{telle} \uncertain{topologie} \uncertain{est} \uncertain{séparée}, \add{\uncertain{\ill{}}}
et \uncertain{a} \uncertain{totalement} \uncertain{discontinue} [\uncertain{pour} \uncertain{deux} \uncertain{pts} $x\neq y$, \uncertain{il} \uncertain{existe} \uncertain{des} \uncertain{ouverts} $V\ni x$, $W\ni y$, $V\cap W=\emptyset$, $V\cup W=U$].
L'hyp.\ (ii) exprime \uncertain{qu'elle} \uncertain{est} \uncertain{quasi-compacte}, \uncertain{donc} \emph{compacte} [\uncertain{dans} \uncertain{l'espace} \ldots{}] \uncertain{totalement} \uncertain{discontinu} \uncertain{et} \ldots, \uncertain{\ill{}}].
\note{la plus grande partie des liaisons de cette page biffée est conjecturale ; seuls les énoncés (i), (ii) et la condition de séparation $V\cap W=\emptyset$, $V\cup W=U$ sont lus avec quelque assurance}

\page{117}
\note{plan numéroté, écrit sur le verso d'une feuille dactylographiée sans rapport (le texte tapé transparaît, à l'envers)}
\begin{itemize}
\item[\textbf{N° 1}] \uncertain{Préschémas} \uncertain{compacts} et loc.\ \uncertain{compacts}. Anneaux flasques.
\item[\textbf{N° 2}] \struck{Systèmes projectifs} \add{Proschémas} flasques \struck{\uncertain{de} préschémas et} \struck{\ill{} \ill{}} \ill{} espaces schématiques \uncertain{compacts}
\item[\textbf{N° 3}] \struck{Espaces \uncertain{schématiques} \ill{}} \struck{\uncertain{Ind} \ill{}} Espaces schématiques. [Catégorie \uncertain{\ill{}} \uncertain{par} les \uncertain{\ill{}} \uncertain{ouverts}\ldots]
\item[\textbf{N° 4}] Opérations \uncertain{sur} \uncertain{ces}, \struck{\ill{}} et \uncertain{lim} \struck{\ill{}} des \uncertain{espaces} \uncertain{schématiques} \add{\uncertain{\ill{}} \uncertain{des} \uncertain{proschémas} $\simeq$ \uncertain{\ill{}} \uncertain{des} \uncertain{compacts}}
\item[\textbf{N° 5}] \struck{\ill{}} Espaces quasi-schématiques\ldots
\end{itemize}
\marginal{Morphismes quasi-compacts, \uncertain{opérations} \uncertain{propres}}\note{note de l'auteur dans la marge gauche, reliée par un trait au N° 3--4}

\note{entre N° 5 et N° 6, un premier diagramme est biffé de boucles et de traits ; on y lit : (Préschémas quasi-compacts quasi-séparés) $\xrightarrow{\alpha}$ (Ens.\ schématiques compacts) $\xrightarrow{\text{épais.}}$ (Pro(Sch) \uncertain{limites} quasi-compacts) $\xrightarrow{\text{fid., \uncertain{surjectif}}}$ (Schémas compacts réduits), avec une flèche de retour « pl.\ fid.\ $\delta$ » et une autre marquée « pl.\ fid. »}

\textbf{N° 6} \uncertain{Pb} \uncertain{d'espaces}. \uncertain{Espaces} \uncertain{de} Néron. \struck{\ill{}}

\begin{tikzcd}[column sep=large, row sep=large, nodes={font=\scriptsize}]
X \arrow[r, "\alpha"] \arrow[rr, bend left=30, "{\gamma\ [\mathrm{Cons\ cpct}]}"] & Y \arrow[r, "\beta=\varprojlim"] & K \arrow[l, bend left=35, "{\beta'\ \text{pl. fid.}}"] \arrow[ll, bend left=55, "{\gamma'\ \text{pl. fid.}}"]
\end{tikzcd}

\note{sur la page, $X$ désigne la catégorie (Préschémas quasi-compacts quasi-séparés), $Y$ la catégorie ($\mathrm{Pro(Sch)}$ flasques), $K$ la catégorie (schémas compacts réduits) ; les lettres sont écrites sous les noms des catégories. L'étiquette de $\gamma'$ se poursuit par un mot lu « \uncertain{immersion} »}
\note{sous la flèche $\alpha$ : « \uncertain{subjectif} [fidèle \uncertain{sur} \uncertain{les} réduits] » ; sous $\beta$ : « \uncertain{surjectif} \uncertain{fidèle} ». Les flèches $\beta'$ et $\gamma'$ reviennent de $K$ en arcs sous la ligne ; $\gamma$ part de $X$ en arc au-dessus. En marge à droite : « $\rho$ [pleinement fidèle, \ill{} \ill{} \ill{} \ill{} \ill{}\ldots] »}

\[
\mathrm{Hom}(K,\gamma X)\xrightarrow{\ \sim\ }\mathrm{Hom}(\gamma' K,X)
\]
\[
\mathrm{Hom}(K,\beta Y)\xrightarrow{\ \sim\ }\mathrm{Hom}(\beta' K,Y)
\]
\[
\mathrm{Hom}(K,\beta(\alpha X))=\mathrm{Hom}(\beta' K,\alpha X)
\]
$\beta'(K)=$ \uncertain{système} \uncertain{projectif} \uncertain{des} \uncertain{\ill{}} \uncertain{définis} \uncertain{par} $K$.
\note{les flèches $\gamma'$ et $\beta'$ de ces formules sont écrites $\gamma^{\cdot}$, $\beta^{\cdot}$ avec un point ou un accent ; lecture du signe incertaine}

\page{119}
\begin{tikzcd}[column sep=large, row sep=large, nodes={font=\scriptsize}]
(\text{Préschémas quasi-séparés}) \arrow[r, "\overline{\alpha}"] \arrow[rr, bend left=30, "\overline{\gamma}"] & \mathrm{Ind}(\mathrm{Cons\,cpct})\text{ flasques} \arrow[r, "\overline{\beta}", "\text{surjectif}"'] & (\text{schémas loc.\ cpcts réduits}) \arrow[l, bend left=35, "{\overline{\beta}'\ \text{pl. fid.}}"] \arrow[ll, bend left=50, "{\overline{\gamma}'\ \text{pl. fid.}}"]
\end{tikzcd}

\note{l'arc $\overline{\gamma}$, au-dessus de la ligne, n'a pas de pointe visible ; son départ (au-dessus de $\mathrm{Ind}$ ou de la première parenthèse) est incertain. Sous $\mathrm{Ind}(\mathrm{Cons\,cpct})$, un « (Cons) » surmonté d'un tréma, lecture incertaine ; « surjectif » sous $\overline{\beta}$ est une lecture incertaine aussi}

\begin{align*}
\mathrm{Hom}(\overline{K},\overline{\beta}\,\overline{Y}) &\simeq \mathrm{Hom}(\overline{\beta}'(K),\overline{Y})
 \simeq \varprojlim_{\lambda}\varinjlim_{\mu}\mathrm{Hom}(\overline{K}_{\lambda},\beta(\overline{Y}_{\mu}))\\
&\simeq \varprojlim_{\lambda}\varinjlim_{\mu}\mathrm{Hom}(\beta'(\overline{K}_{\lambda}),\overline{Y}_{\mu})
\end{align*}
\[
\text{\struck{$\mathrm{Hom}(\overline{K},\overline{\gamma}(X))\simeq$}}
\]
\begin{align*}
\mathrm{Hom}(K,\overline{\gamma}(X)) &\simeq \mathrm{Hom}(\overline{\gamma}'K,X)\\
&\simeq \varprojlim_{\lambda}\varinjlim_{\mu}\mathrm{Hom}(K_{\lambda},\gamma(X_{\mu}))\\
&\simeq \varprojlim_{\lambda}\varinjlim_{\mu}\mathrm{Hom}(\gamma' K_{\lambda},X_{\mu})
\end{align*}
\note{un trait courbe relie la deuxième ligne de droite ($\beta'(\overline{K}_{\lambda})$) à la formule en $\overline{\gamma}$ ; dans la dernière ligne, un caractère est écrasé devant $X_{\mu}$ ; le premier $\overline{K}$ porte une double barre, peut-être une flèche}

\begin{tikzcd}[column sep=large]
(\mathrm{Cons}) \arrow[r, "\text{pl. fid.}"] & \underline{\mathrm{Hom}}((\mathrm{Cons\,cpct})^{\circ},(\mathrm{Ens}))
\end{tikzcd}

\note{fin de la page et du lot ; le développement sur les espaces schématiques (« Ind » d'ensembles constructibles compacts) se poursuit vraisemblablement au-delà de la page 120}

\end{document}
